\chapter{Teoría de grupos}\label{ch:b3:groups}

En el volumen de segundo año los grupos eran poco más que un instrumento
de recuento: teorema de Lagrange, grupos cíclicos, grupo simétrico y su
signatura. Este capítulo convierte la teoría de grupos en un
\emph{método}. El motor es la noción de grupo que \emph{actúa} sobre un
conjunto: contar \omterm{def:b3:groups:action}{órbitas} y puntos fijos
produce la \omterm{cor:b3:groups:classeq}{ecuación de clases}, el
teorema de Cauchy y los tres teoremas de Sylow --- el principio
local-global fundamental de la teoría de grupos finitos. Aprenderemos
después a ensamblar grupos (\omterm{prop:b3:groups:direct}{productos directos} y
semidirectos) y a desmontarlos
(\omterm{thm:b3:groups:jordanholder}{series de composición},
\omterm{def:b3:groups:derived}{grupos resolubles}), y demostraremos el
teorema que, en el~\cref{ch:b3:galois}, cerrará una cuestión de tres
siglos sobre las ecuaciones polinómicas: el
\omterm{thm:b3:groups:ansimple}{grupo alternado} $A_n$ es
\omterm{def:b3:groups:simple}{simple} para $n
\geq 5$.

\section{Grupos cociente y teoremas de isomorfía}

En todo el capítulo, $G$ es un grupo escrito multiplicativamente y $e$
su elemento neutro. Recordemos del volumen de segundo año: los
subgrupos, las clases laterales $gH$, el teorema de Lagrange
($\abs G = [G:H]\,\abs H$ para $G$ finito), el orden de un elemento, los
grupos cíclicos y el grupo simétrico $S_n$ con su morfismo signatura
$\varepsilon \colon S_n \to \{\pm1\}$.

\begin{definition}\label{def:b3:groups:normal}
Un subgrupo $N$ de $G$ es \emph{normal}\index{subgrupo normal} (se
escribe $N \trianglelefteq G$) cuando $gNg^{-1} = N$ para todo $g
\in G$ --- equivalentemente, cuando las clases laterales por la
izquierda y por la derecha coinciden: $gN =
Ng$ para todo $g$.
\end{definition}

\begin{theorem}[Grupo cociente]\label{thm:b3:groups:quotient}
Sea $N \trianglelefteq G$. El conjunto $G/N$ de las clases laterales,
dotado de la multiplicación $(gN)(hN) = ghN$, es un grupo bien definido,
el \emph{grupo cociente}\index{grupo cociente}, y la \emph{proyección
canónica} $\pi \colon G \to G/N$, $g \mapsto gN$, es un morfismo
sobreyectivo de núcleo $N$. Recíprocamente, todo núcleo de un morfismo
de grupos es \omterm{def:b3:groups:normal}{normal}: los
\omterm{def:b3:groups:normal}{subgrupos normales} son exactamente los
núcleos.
\end{theorem}

\begin{proof}
Lo esencial es la buena definición. Si $gN = g'N$ y $hN = h'N$,
escribamos $g' = gn$, $h' = hm$ con $n, m \in N$. Entonces
$g'h' = gnhm =
gh\,(h^{-1}nh)\,m \in ghN$, pues $h^{-1}nh \in N$ por normalidad: el
producto de clases no depende de los representantes. La asociatividad,
el neutro $eN = N$ y los inversos $(gN)^{-1} = g^{-1}N$ se heredan de
$G$. Es claro que $\pi$ es un morfismo sobreyectivo y
$\pi(g) = N \iff g \in N$.

Si $f \colon G \to H$ es un morfismo y $k \in \ker f$, entonces
$f(gkg^{-1}) = f(g)f(k)f(g)^{-1} = e$: los núcleos son
\omterm{def:b3:groups:normal}{normales}.
\end{proof}

\begin{theorem}[Propiedad universal; primer teorema de isomorfía]
\label{thm:b3:groups:firstiso}
Sea $f \colon G \to H$ un morfismo y sea $N \trianglelefteq G$ con
$N \subseteq \ker f$. Existe un único morfismo $\bar f \colon G/N
\to H$ tal que $f = \bar f \circ \pi$. En particular, tomando $N = \ker
f$:\index{teoremas de isomorfía}
\[
G/\ker f \;\xrightarrow{\;\sim\;}\; \operatorname{im} f,
\qquad gN \mapsto f(g).
\]
\end{theorem}

\begin{proof}
Unicidad: $\bar f(gN)$ ha de ser $f(g)$. Existencia: si $gN = g'N$,
entonces $g^{-1}g' \in N \subseteq \ker f$, luego $f(g) = f(g')$ y $\bar
f(gN) = f(g)$ está bien definida; es un morfismo porque lo es $f$. Para
$N = \ker f$: $\bar f$ es inyectiva, ya que $\bar f(gN) = e$ significa
$g \in \ker f$, es decir, $gN = N$; su imagen es la de $f$.
\end{proof}

\begin{theorem}[Segundo y tercer teoremas de isomorfía]
\label{thm:b3:groups:secondiso}
Sean $H \leq G$ y $N \trianglelefteq G$.
\begin{enumerate}
  \item $HN = \{hn : h \in H,\, n \in N\}$ es un subgrupo, $N
        \trianglelefteq HN$, $H \cap N \trianglelefteq H$, y
        \[
        H/(H \cap N) \;\cong\; HN/N .
        \]
  \item Si además $N \subseteq K \trianglelefteq G$, entonces $K/N
        \trianglelefteq G/N$ y $(G/N)\big/(K/N) \cong G/K$.
\end{enumerate}
\end{theorem}

\begin{proof}
(1) $HN$ es un subgrupo: $(hn)(h'n') = hh'\,(h'^{-1}nh')n' \in HN$ y
$(hn)^{-1} = h^{-1}(hn^{-1}h^{-1}) \in HN$, usando la normalidad de $N$.
Compongamos $H \hookrightarrow HN \xrightarrow{\pi} HN/N$: este morfismo
es sobreyectivo ($hnN = hN$) y tiene núcleo $\{h \in H : h \in
N\} = H \cap N$; basta aplicar el~\cref{thm:b3:groups:firstiso}.

(2) La proyección $G/N \to G/K$, $gN \mapsto gK$, está bien definida
($N \subseteq K$), es sobreyectiva y tiene núcleo $K/N$; se aplica de
nuevo el~\cref{thm:b3:groups:firstiso}.
\end{proof}

\begin{theorem}[Teorema de
correspondencia]\label{thm:b3:groups:correspondence}
Sea $N \trianglelefteq G$. La aplicación $H \mapsto H/N$ es una
biyección entre los subgrupos de $G$ que contienen $N$ y los subgrupos
de $G/N$, y conserva inclusiones, índices y normalidad (en ambos
sentidos).
\end{theorem}

\begin{proof}
Su inversa es $\bar H \mapsto \pi^{-1}(\bar H)$. Ambas aplicaciones
llevan subgrupos en subgrupos y son mutuamente inversas:
$\pi^{-1}(H/N) =
HN = H$ porque $N \subseteq H$, y $\pi(\pi^{-1}(\bar H)) = \bar H$ por
la sobreyectividad de $\pi$. Las inclusiones se conservan claramente;
$[G:H] = [G/N : H/N]$ porque $gH \mapsto (gN)(H/N)$ es una biyección
bien definida entre los espacios de clases; y $gHg^{-1} = H$ para todo
$g$ si y solo si $(gN)(H/N)(gN)^{-1} = H/N$ para todo $gN$, de nuevo por
la sobreyectividad de $\pi$.
\end{proof}

\begin{example}\label{ex:b3:groups:firstiso}
$\varepsilon \colon S_n \to \{\pm 1\}$ da $S_n/A_n \cong
\{\pm1\}$; $\det \colon GL_n(K) \to K^\times$ da
$GL_n(K)/SL_n(K) \cong K^\times$; $t \mapsto \eu^{2\iu\pi t}$ da
$\R/\Z \cong \mathbb U$, el grupo del círculo. El primer teorema de
isomorfía es la forma en que los cocientes se \emph{calculan} en la
práctica: buscar una sobreyección con el núcleo adecuado.
\end{example}

\begin{method}\label{met:b3:groups:normal}
Para demostrar que $N \trianglelefteq G$, en orden decreciente de
elegancia: exhibir $N$ como núcleo de un morfismo definido sobre $G$;
comprobar $gNg^{-1} \subseteq N$ para todo $g$ (basta con eso:
aplicándolo a $g^{-1}$ y conjugando se obtiene la inclusión recíproca);
verificar que $N$ es una unión de \omterm{ex:b3:groups:actions}{clases
de conjugación}; o bien observar que $[G:N] = 2$ (entonces $gN = Ng$ es
forzoso --- el~\cref{exo:b3:groups:1}).
\end{method}

\section{Acciones de grupo}

\begin{definition}\label{def:b3:groups:action}
Una \emph{acción}\index{acción de grupo} de $G$ sobre un conjunto $X$ es
un morfismo $\varphi \colon G \to \mathfrak{S}(X)$ en el grupo de las
biyecciones de $X$; se escribe $g \cdot x$ en lugar de $\varphi(g)(x)$.
Equivalentemente: una aplicación $G \times X \to X$ con $e \cdot x = x$
y $g
\cdot (h \cdot x) = (gh) \cdot x$. La \emph{órbita}\index{órbita} de $x$
es $\mathcal O_x = \{g \cdot x : g \in G\}$, su
\emph{estabilizador}\index{estabilizador} es el subgrupo $G_x = \{g :
g\cdot x = x\}$, y $X^G = \{x : \forall g,\ g \cdot x = x\}$ es el
conjunto de los \emph{puntos fijos}. La acción es \emph{transitiva} si
hay exactamente una órbita, \emph{fiel} si $\varphi$ es inyectivo y
\emph{libre} si todos los estabilizadores son triviales.
\end{definition}

\begin{example}\label{ex:b3:groups:actions}
Cinco acciones gobiernan toda la teoría
de grupos finitos:
\begin{enumerate}
  \item $G$ sobre sí mismo por \emph{traslación por la izquierda}
  $g \cdot x = gx$: libre y transitiva. \item $G$ sobre sí mismo por
  \emph{conjugación} $g \cdot x = gxg^{-1}$: las
  \omterm{def:b3:groups:action}{órbitas} son las \emph{clases de
  conjugación}\index{clase de conjugación}, los
  \omterm{def:b3:groups:action}{estabilizadores} los
  \emph{centralizadores} \index{centralizador}
  $Z_G(x) = \{g : gx = xg\}$ y los puntos fijos el
  \emph{centro}\index{centro} $Z(G)$. \item $G$ sobre el espacio de
  clases $G/H$ por $g \cdot xH = gxH$: transitiva, con
  \omterm{def:b3:groups:action}{estabilizador} de la clase $H$ igual a
  $H$. Toda \omterm{def:b3:groups:action}{acción} transitiva es de esta
  forma
        (\cref{exo:b3:groups:8}).
  \item $G$ sobre el conjunto de sus subgrupos por conjugación: el
  \omterm{def:b3:groups:action}{estabilizador} de $H$ es el
  \emph{normalizador}\index{normalizador} $N_G(H) =
        \{g : gHg^{-1} = H\}$, el mayor subgrupo de $G$ en el que $H$ es
  \omterm{def:b3:groups:normal}{normal}. \item $S_n$ sobre
  $\intint{1}{n}$: la madre de todos los ejemplos.
\end{enumerate}
\end{example}

\begin{theorem}[Órbita--estabilizador]\label{thm:b3:groups:orbitstab}
La aplicación $gG_x \mapsto g \cdot x$ es una biyección bien definida
$G/G_x \to \mathcal O_x$. En particular, para $G$ finito,\index{teorema
órbita--estabilizador}
\[
\abs{\mathcal O_x} = [G : G_x] \quad\text{divide a } \abs G,
\]
y, como las \omterm{def:b3:groups:action}{órbitas} forman una partición
de $X$ (son las clases de la equivalencia
$x \sim y \iff y \in \mathcal O_x$),
\[
\abs X = \sum_{i} \,[G : G_{x_i}]
\qquad (x_i\colon \text{un punto por órbita}).
\]
\end{theorem}

\begin{proof}
Buena definición e inyectividad: $gG_x = hG_x \iff h^{-1}g \in G_x \iff
h^{-1}g \cdot x = x \iff g \cdot x = h \cdot x$; léase la cadena en
ambos sentidos. La sobreyectividad es la definición misma de
\omterm{def:b3:groups:action}{órbita}. Los enunciados de recuento se
siguen del teorema de Lagrange y de la partición de $X$ en
\omterm{def:b3:groups:action}{órbitas}.
\end{proof}

\begin{corollary}[Ecuación de clases]\label{cor:b3:groups:classeq}
Para un grupo finito $G$, eligiendo un representante $x_i$ en cada
\omterm{ex:b3:groups:actions}{clase de conjugación} con más de un
elemento:\index{ecuación de clases}
\[
\abs G = \abs{Z(G)} + \sum_i \,[G : Z_G(x_i)],
\qquad\text{cada } [G : Z_G(x_i)] > 1 \text{ que divide a } \abs G.
\]
\end{corollary}

\begin{proof}
Basta aplicar el~\cref{thm:b3:groups:orbitstab} a la
\omterm{def:b3:groups:action}{acción} por conjugación: las
\omterm{def:b3:groups:action}{órbitas} con un solo elemento son
exactamente los elementos de $Z(G)$.
\end{proof}

\begin{theorem}[Puntos fijos de los
$p$-grupos]\label{thm:b3:groups:pfixed}
Sea $p$ primo. Un \emph{$p$-grupo}\index{p-grupo@$p$-grupo} es un grupo
finito cuyo orden es una potencia de $p$. Si un $p$-grupo $G$ actúa
sobre un conjunto finito $X$, entonces
\[
\abs{X^G} \equiv \abs X \pmod p .
\]
Consecuencias: un $p$-grupo no trivial tiene
\omterm{ex:b3:groups:actions}{centro} no trivial, y todo grupo de orden
$p^2$ es abeliano.
\end{theorem}

\begin{proof}
Cada \omterm{def:b3:groups:action}{órbita} tiene cardinal $[G:G_x]$, una
potencia de $p$; esa potencia vale $1$ exactamente sobre los puntos
fijos y en los demás casos es divisible por $p$. Sumando sobre las
\omterm{def:b3:groups:action}{órbitas} se obtiene la congruencia. Para
el \omterm{ex:b3:groups:actions}{centro}: la
\omterm{def:b3:groups:action}{acción} de $G$ sobre sí mismo por
conjugación cumple $X^G = Z(G)$, luego $\abs{Z(G)} \equiv
\abs G \equiv 0 \pmod p$, y $Z(G) \ni e$ obliga a $\abs{Z(G)} \geq
p$. Orden $p^2$: si $Z(G) \neq G$, entonces $\abs{Z(G)} = p$ y $G/Z(G)$
es cíclico de orden $p$, lo que obliga a que $G$ sea abeliano
(\cref{exo:b3:groups:2}) --- contradicción.
\end{proof}

\begin{theorem}[Cauchy]\label{thm:b3:groups:cauchy}
Si un primo $p$ divide a $\abs G$, entonces $G$ contiene un elemento de
orden $p$.\index{teorema de Cauchy}
\end{theorem}

\begin{proof}[Demostración (McKay)]
Sea $X = \{(g_1, \dots, g_p) \in G^p : g_1 g_2 \cdots g_p = e\}$. Al
elegir libremente $g_1, \dots, g_{p-1}$ queda determinado $g_p$:
$\abs X =
\abs G^{p-1}$, divisible por $p$. El grupo cíclico $\Z/p\Z$ actúa sobre
$X$ por desplazamiento cíclico
$(g_1, \dots, g_p) \mapsto (g_2, \dots, g_p,
g_1)$ --- esto preserva $X$, pues $g_2 \cdots g_p g_1 =
g_1^{-1}(g_1 \cdots g_p)g_1 = e$. Por el~\cref{thm:b3:groups:pfixed},
$\abs{X^{\Z/p\Z}} \equiv \abs X \equiv 0 \pmod p$. Los puntos fijos son
las tuplas constantes $(g, \dots, g)$ con $g^p = e$; la tupla $(e,
\dots, e)$ es uno de ellos, de modo que hay al menos $p$, y por tanto al
menos un $g \neq e$ con $g^p = e$: su orden es exactamente $p$.
\end{proof}

\begin{theorem}[Cayley]\label{thm:b3:groups:cayley}
Todo grupo de orden $n$ se sumerge en $S_n$.\index{teorema de Cayley}
\end{theorem}

\begin{proof}
La traslación por la izquierda
$\varphi \colon G \to \mathfrak S(G) \cong S_n$ es un morfismo;
$\varphi(g) = \mathrm{id}$ obliga a $g = ge = e$: es fiel.
\end{proof}

\begin{method}\label{met:b3:groups:counting}
El recuento de puntos fijos es la apertura universal de la teoría de
grupos finitos. Para demostrar que algo \emph{existe} (un elemento
central, un elemento de orden $p$, un
\omterm{def:b3:groups:normal}{subgrupo normal}, un punto fijo), se hace
actuar un grupo bien elegido sobre un conjunto finito bien elegido y se
comparan $\abs{X^G}$ y $\abs X$ módulo $p$, o bien se obliga a los
tamaños de las \omterm{def:b3:groups:action}{órbitas} a dividir el orden
del grupo. Las demostraciones
del~\cref{thm:b3:groups:pfixed,thm:b3:groups:cauchy} y de los tres
teoremas de Sylow que siguen son cinco variaciones sobre esta única
idea.
\end{method}

\section{Los teoremas de Sylow}

El teorema de Lagrange afirma que el orden de un subgrupo divide a
$\abs G$; el recíproco es falso ($A_4$, de orden $12$, no tiene ningún
subgrupo de orden $6$ --- el~\cref{exo:b3:groups:1}). Los teoremas de
Sylow rescatan el recíproco para las potencias de primos, y su cláusula
de recuento es la herramienta general más fina de que disponemos para
producir \omterm{def:b3:groups:normal}{subgrupos normales}.

\begin{definition}\label{def:b3:groups:sylow}
Escribamos $\abs G = p^a m$ con $p \nmid m$. Un \emph{$p$-subgrupo de
Sylow}\index{subgrupo de Sylow} de $G$ es un subgrupo de orden $p^a$ ---
un $p$-subgrupo del mayor orden concebible. El número de $p$-subgrupos
de Sylow de $G$ se denota $n_p$.
\end{definition}

\begin{lemma}\label{lem:b3:groups:binom}
Si $\abs G = p^a m$ con $p \nmid m$, entonces $\dbinom{p^a m}{p^a}
\equiv m \pmod p$.
\end{lemma}

\begin{proof}
En $\mathbb F_p[X]$, el sueño del principiante $(1+X)^p = 1 + X^p$ (los
coeficientes $\binom pk$, $0<k<p$, son divisibles por $p$: $p$ divide al
numerador de $\frac{p!}{k!(p-k)!}$ pero no al denominador) se itera
hasta $(1+X)^{p^a} = 1 + X^{p^a}$, de donde
\[
(1+X)^{p^a m} = \bigl(1 + X^{p^a}\bigr)^m
= \sum_{k=0}^{m} \binom mk X^{k p^a} \quad\text{en } \mathbb
F_p[X].
\]
Identifíquese el coeficiente de $X^{p^a}$: a la izquierda $\binom{p^a
m}{p^a} \bmod p$, a la derecha $\binom m1 = m$.
\end{proof}

\begin{theorem}[Sylow I: existencia]\label{thm:b3:groups:sylow1}
Para todo primo $p$ existen $p$-subgrupos de Sylow de $G$.
\end{theorem}

\begin{proof}[Demostración (Wielandt)]
Sea $\Omega$ el conjunto de los \emph{subconjuntos} de $G$ de cardinal
$p^a$; $G$ actúa sobre $\Omega$ por traslación por la izquierda
$g \cdot S = gS$. Por el~\cref{lem:b3:groups:binom},
$\abs\Omega = \binom{p^a m}{p^a}
\equiv m \not\equiv 0 \pmod p$, luego alguna
\omterm{def:b3:groups:action}{órbita} $\mathcal O_S$ tiene tamaño primo
con $p$ (si $p$ dividiese el tamaño de todas las
\omterm{def:b3:groups:action}{órbitas}, dividiría a $\abs\Omega$). Sea
$H = G_S$ el \omterm{def:b3:groups:action}{estabilizador} de uno de esos
$S$. Como $[G : H] = \abs{\mathcal O_S}$ es primo con $p$ y $p^a \mid
\abs G = [G:H]\,\abs H$, obtenemos $p^a \mid \abs H$. Recíprocamente,
fijemos $s \in S$: la aplicación $H \to S$, $h \mapsto hs$, es inyectiva
y toma valores en $S$, porque $hS = S$; por tanto
$\abs H \leq \abs S = p^a$. Así pues, $\abs H = p^a$.
\end{proof}

\begin{theorem}[Sylow II: dominación y
conjugación]\label{thm:b3:groups:sylow2}
Sea $P$ un $p$-subgrupo de Sylow y sea $Q$ un $p$-subgrupo cualquiera de
$G$. Entonces $Q \subseteq gPg^{-1}$ para algún $g \in G$. En
particular, todos los $p$-subgrupos de Sylow son conjugados, y
$P \trianglelefteq G \iff
n_p = 1$.
\end{theorem}

\begin{proof}
Hagamos actuar $Q$ sobre el espacio de clases $X = G/P$, de cardinal $m
\not\equiv 0 \pmod p$. Por el~\cref{thm:b3:groups:pfixed} aplicado al
$p$-grupo $Q$, $\abs{X^Q} \equiv m \not\equiv 0 \pmod p$: existe una
clase fija $gP$, es decir, $QgP = gP$, o sea $g^{-1}Qg \subseteq
P$. Si $Q$ es a su vez un \omterm{def:b3:groups:sylow}{subgrupo de
Sylow}, la igualdad de órdenes convierte $Q \subseteq gPg^{-1}$ en una
igualdad. Por último, $P
\trianglelefteq G$ si y solo si sus conjugados $\{gPg^{-1}\}$ --- que,
por lo anterior, son \emph{todos} los $p$-subgrupos de Sylow --- se
reducen a
$\{P\}$.
\end{proof}

\begin{theorem}[Sylow III: recuento]\label{thm:b3:groups:sylow3}
$n_p \equiv 1 \pmod p$, y $n_p = [G : N_G(P)]$, que divide a
$m$.
\end{theorem}

\begin{proof}
Sea $\mathrm{Syl}_p$ el conjunto de los $p$-subgrupos de Sylow; $G$
actúa sobre él transitivamente por conjugación
(\cref{thm:b3:groups:sylow2}), con
\omterm{def:b3:groups:action}{estabilizador} de $P$ igual al
\omterm{ex:b3:groups:actions}{normalizador} $N_G(P) \supseteq P$:
$n_p = [G :
N_G(P)]$, y $m = [G:P] = [G:N_G(P)]\,[N_G(P):P]$ muestra que $n_p \mid
m$.

Restrinjamos ahora la \omterm{def:b3:groups:action}{acción} a $P$ y
contemos los puntos fijos. Si $Q \in
\mathrm{Syl}_p$ queda fijo por $P$, entonces $P \subseteq N_G(Q)$; tanto
$P$ como $Q$ son $p$-subgrupos de Sylow del grupo $N_G(Q)$, luego son
conjugados en él (el~\cref{thm:b3:groups:sylow2} aplicado a $N_G(Q)$);
pero $Q \trianglelefteq N_G(Q)$, de modo que $Q$ es su único conjugado
allí: $P = Q$. Así, el único punto fijo es el propio $P$, y
el~\cref{thm:b3:groups:pfixed} da $n_p = \abs{\mathrm{Syl}_p}
\equiv \abs{\mathrm{Syl}_p^P} = 1 \pmod p$.
\end{proof}

\begin{method}\label{met:b3:groups:sylowmethod}
Para analizar un grupo de orden dado $n = p^a m$: listar los divisores
de $m$ congruentes con $1$ módulo $p$ --- son los candidatos a $n_p$. Si
el único candidato es $1$, el $p$-subgrupo de Sylow es
\omterm{def:b3:groups:normal}{normal}. Si $n_p > 1$ ha de ser pequeño,
hágase actuar por conjugación sobre $\mathrm{Syl}_p$ para obtener un
morfismo $G \to S_{n_p}$ de núcleo pequeño. Y cuéntense elementos: dos
$p$-subgrupos de Sylow distintos de orden \emph{primo} $p$ se cortan
trivialmente, así que aportan $n_p(p-1)$ elementos de orden exactamente
$p$; los recuentos para primos distintos se solapan y a menudo fuerzan
una contradicción
(\cref{exo:b3:groups:7}).
\end{method}

\begin{example}\label{ex:b3:groups:pq}
Sea $\abs G = pq$ con $p < q$ primos y $p \nmid q - 1$. Entonces
$n_q \mid p$ y $n_q \equiv 1 \bmod q$ obligan a $n_q = 1$ (pues $p <
q$); $n_p \mid q$ y $n_p \equiv 1 \bmod p$ obligan a $n_p = 1$ (pues $q
\not\equiv 1 \bmod p$). Sean $P, Q$ los dos \omterm{def:b3:groups:sylow}{subgrupos de Sylow}
\omterm{def:b3:groups:normal}{normales}: $P
\cap Q = \{e\}$ (órdenes coprimos), luego $\abs{PQ} = pq$
(\cref{exo:b3:groups:4}) y $G \cong P \times Q \cong \Z/p\Z
\times \Z/q\Z \cong \Z/pq\Z$ por la~\cref{prop:b3:groups:direct} de más
abajo. Todo grupo de orden $15$, $33$, $35$, \dots\ es cíclico. El caso
excluido $p \mid q - 1$ produce exactamente un grupo más, no abeliano
--- véase el problema de fin de semana (\cref{pb:b3:groups:1}).
\end{example}

\begin{example}[Un censo de Sylow completo: $S_4$]
\label{ex:b3:groups:s4sylow}
Apliquemos el método a $G = S_4$, $\abs G = 24 = 2^3\cdot3$. \emph{Sylow
$3$:} $n_3 \mid 8$, $n_3 \equiv 1 \bmod 3$, luego $n_3
\in \{1, 4\}$; como $\langle(123)\rangle$ y $\langle(124)\rangle$ son
distintos, $n_3 = 4$ --- los cuatro subgrupos $\langle(abc)\rangle$, uno
por cada subconjunto $3$ de $\{a, b, c\}$ elementos, que dan cuenta de
los $8$ ciclos de longitud tres. Por Sylow~II son conjugados, y aquí el
morfismo de conjugación $S_4 \to
S_{\mathrm{Syl}_3} \cong S_4$ es un isomorfismo (su núcleo está
contenido en $N = N_G(\langle(123)\rangle)$, de orden $24/4 =
6$, y un \omterm{def:b3:groups:normal}{subgrupo normal} de $S_4$ dentro
de un $S_3$ como $N$ ha de ser trivial: constaría de permutaciones pares
que fijan los cuatro \omterm{def:b3:groups:sylow}{subgrupos de Sylow}, y solo $e$ lo hace).
\emph{Sylow $2$:} $n_2 \mid 3$, $n_2
\equiv 1 \bmod 2$: $n_2 \in \{1, 3\}$. El subgrupo $D =
\langle(1234), (13)\rangle$ tiene orden $8$ (un diédrico $D_4$: las
simetrías del cuadrado de vértices $1, 2, 3, 4$), no es
\omterm{def:b3:groups:normal}{normal} ($(12)(1234)(12) = (2134)$ genera
un subgrupo cíclico de orden $4$ distinto), luego $n_2 = 3$: las tres
copias de $D_4$ corresponden a las tres maneras de emparejar $4$ puntos
formando un «cuadrado». Obsérvese la moraleja del censo:
$\abs{S_4} = 24$ deja margen para que cualquiera de los dos Sylow deje
de ser \omterm{def:b3:groups:normal}{normal}, y ambos lo hacen --- compárese con el orden $12$, donde
el recuento obliga a que uno de ellos sea
\omterm{def:b3:groups:normal}{normal} (Parte~IV
del~\cref{pb:b3:groups:1}).
\end{example}

\section{Productos, directos y semidirectos}

\begin{proposition}[Reconocer un producto
directo]\label{prop:b3:groups:direct}
Sean $H, K \trianglelefteq G$ con $H \cap K = \{e\}$ y $HK = G$.
Entonces $(h,k) \mapsto hk$ es un isomorfismo $H \times K \to
G$.\index{producto directo}
\end{proposition}

\begin{proof}
Para $h \in H$, $k \in K$, el \omterm{def:b3:groups:derived}{conmutador}
$hkh^{-1}k^{-1}$ está en $K$ (léase como $(hkh^{-1})k^{-1}$, usando la
normalidad de $K$) y en $H$ (léase como $h(kh^{-1}k^{-1})$): vale $e$,
luego $H$ y $K$ conmutan elemento a elemento y la aplicación es un
morfismo. Es sobreyectiva porque $HK = G$, e inyectiva porque $hk = e$
da $h = k^{-1} \in
H \cap K = \{e\}$.
\end{proof}

Lo que más veces falla es la normalidad de \emph{ambos} factores: en
$S_3
= \langle (1\,2\,3)\rangle \,\langle(1\,2)\rangle$ los dos factores se
cortan trivialmente y generan el grupo, y sin embargo
$S_3 \not\cong \Z/3\Z \times
\Z/2\Z$. La noción adecuada cuando solo uno de los factores es
\omterm{def:b3:groups:normal}{normal} es la siguiente:

\begin{definition}\label{def:b3:groups:semidirect}
Sean $H$, $K$ grupos y $\varphi \colon K \to
\operatorname{Aut}(H)$ un morfismo. El \emph{producto
semidirecto}\index{producto semidirecto} $H \rtimes_\varphi K$ es el
conjunto $H \times K$ dotado de
\[
(h, k)\,(h', k') = \bigl(h\,\varphi(k)(h'),\; kk'\bigr).
\]
\end{definition}

\begin{proposition}\label{prop:b3:groups:semidirect}
$H \rtimes_\varphi K$ es un grupo; $H \times \{e\}$ es un
\omterm{def:b3:groups:normal}{subgrupo normal} isomorfo a $H$ y
$\{e\} \times K$ un subgrupo isomorfo a $K$; se cortan trivialmente y
generan el grupo. Recíprocamente, si $G =
NK$ con $N \trianglelefteq G$, $K \leq G$ y $N \cap K = \{e\}$, entonces
$G \cong N \rtimes_\varphi K$ para $\varphi(k) = (n \mapsto
knk^{-1})$.
\end{proposition}

\begin{proof}
Verificación directa: la asociatividad se reduce a $\varphi(kk') =
\varphi(k)\circ\varphi(k')$ y a que cada $\varphi(k)$ sea un morfismo;
el neutro es $(e,e)$ y $(h,k)^{-1} =
\bigl(\varphi(k^{-1})(h^{-1}), k^{-1}\bigr)$. La proyección $(h,k)
\mapsto k$ es un morfismo sobre $K$ de núcleo $H \times \{e\}$, que por
tanto es \omterm{def:b3:groups:normal}{normal}. Para el recíproco: todo
$g \in G$ se escribe de manera \emph{única} como $nk$ con $n \in N$,
$k \in K$ (existencia: $G =
NK$; unicidad: $nk = n'k'$ da $n'^{-1}n = k'k^{-1} \in N \cap
K$), y
\[
(nk)(n'k') = n\,(kn'k^{-1})\;kk'
\]
muestra que $nk \mapsto (n, k)$ transporta la ley de $G$ a la de
$N \rtimes_\varphi K$.
\end{proof}

\begin{example}\label{ex:b3:groups:semidirectexamples}
(a) El \emph{grupo diédrico}\index{grupo diédrico} $D_n$ ($n \geq
3$) de las $2n$ simetrías de un $n$-ágono regular: las rotaciones forman
un \omterm{def:b3:groups:normal}{subgrupo normal} de índice $2$,
cualquier reflexión genera un complemento, y conjugar una rotación por
una reflexión la invierte: $D_n \cong \Z/n\Z \rtimes_\varphi \Z/2\Z$ con
$\varphi(1) = (x
\mapsto -x)$. (b) El grupo afín de una recta,
$\{x \mapsto ax + b : a \in
K^\times,\, b \in K\} \cong K \rtimes K^\times$: las traslaciones son
\omterm{def:b3:groups:normal}{normales} y las homotecias forman un
complemento. (c) $S_n \cong A_n \rtimes \Z/2\Z$ (complemento: cualquier
trasposición). (d) El \omterm{pb:b3:groups:1}{grupo de cuaterniones}
$Q_8$ \emph{no} es \omterm{def:b3:groups:semidirect}{producto
semidirecto} de subgrupos propios: todo subgrupo no trivial contiene
$-1$ (\cref{pb:b3:groups:1}), así que no hay dos subgrupos propios que
se corten trivialmente.
\end{example}

\section{Grupos resolubles; simplicidad de \texorpdfstring{$A_n$}{An}}

\begin{definition}\label{def:b3:groups:derived}
El \emph{conmutador}\index{conmutador} de $x, y \in G$ es $[x,y]
= xyx^{-1}y^{-1}$; el \emph{subgrupo derivado}\index{subgrupo derivado}
$D(G)$ es el subgrupo generado por todos los conmutadores. La
\emph{serie derivada} es $D^0(G) = G$, $D^{i+1}(G) = D(D^i(G))$, y $G$
es \emph{resoluble}\index{grupo resoluble} si $D^n(G) =
\{e\}$ para algún $n$.
\end{definition}

\begin{proposition}\label{prop:b3:groups:derived}
$D(G)$ es \omterm{def:b3:groups:normal}{normal} (de hecho, estable por
todo automorfismo), $G/D(G)$ es abeliano y, para $N \trianglelefteq G$:
$G/N$ abeliano $\iff
D(G) \subseteq N$. Además, $G$ es
\omterm{def:b3:groups:derived}{resoluble} si y solo si existe una cadena
$G = G_0 \trianglerighteq G_1 \trianglerighteq \dots
\trianglerighteq G_n = \{e\}$ con cada $G_{i+1} \trianglelefteq
G_i$ y cada cociente $G_i/G_{i+1}$ abeliano. Los subgrupos y los
cocientes de \omterm{def:b3:groups:derived}{grupos resolubles} son
\omterm{def:b3:groups:derived}{resolubles}; recíprocamente, si $N$ y
$G/N$ son \omterm{def:b3:groups:derived}{resolubles}, también lo es $G$.
\end{proposition}

\begin{proof}
Un automorfismo $\alpha$ lleva $[x,y]$ en $[\alpha x, \alpha y]$:
permuta los \omterm{def:b3:groups:derived}{conmutadores}, luego conserva
el subgrupo que generan; las conjugaciones son automorfismos, de donde
la normalidad. En $G/D(G)$,
$\bar x\bar y \bar x^{-1}\bar y^{-1} = \overline{[x,y]} = \bar e$:
el cociente es abeliano. Si $G/N$ es abeliano, entonces todo $[x,y] \in
N$, luego $D(G) \subseteq N$; recíprocamente, si $D(G) \subseteq N$,
entonces $G/N$, que es un cociente del abeliano $G/D(G)$ por el tercer
teorema de isomorfía, es abeliano.

Si $G$ es \omterm{def:b3:groups:derived}{resoluble}, la
\omterm{def:b3:groups:derived}{serie derivada} es una cadena de este
tipo. Recíprocamente, dada una cadena, $D^i(G) \subseteq G_i$ por
inducción: $G_i/G_{i+1}$ abeliano da $D(G_i) \subseteq G_{i+1}$, luego
$D^{i+1}(G) = D(D^i G)
\subseteq D(G_i) \subseteq G_{i+1}$; por tanto $D^n(G) = \{e\}$.

Herencia: $D^i(H) \subseteq D^i(G)$ para $H \leq G$ (inducción), y
$D^i(G/N) = \pi(D^i(G))$ porque $\pi$ lleva
\omterm{def:b3:groups:derived}{conmutadores} sobre
\omterm{def:b3:groups:derived}{conmutadores}; de ahí los enunciados para
subgrupos y cocientes. Extensión: si $D^m(G/N) = \{e\}$, entonces
$D^m(G) \subseteq N$, y $D^n(N) = \{e\}$ da
$D^{m+n}(G) = D^n(D^m(G)) \subseteq D^n(N)
= \{e\}$.
\end{proof}

\begin{example}\label{ex:b3:groups:solvableexamples}
Los grupos abelianos son \omterm{def:b3:groups:derived}{resolubles}. Los
$p$-grupos son \omterm{def:b3:groups:derived}{resolubles}, por inducción
sobre el orden: $Z(G) \neq \{e\}$ y $G/Z(G)$ es un $p$-grupo más
pequeño. $S_3$ y $S_4$ son \omterm{def:b3:groups:derived}{resolubles}:
$S_4 \trianglerighteq A_4
\trianglerighteq V \trianglerighteq \{e\}$, donde $V = \{e,
(1\,2)(3\,4), (1\,3)(2\,4), (1\,4)(2\,3)\}$ es el grupo de Klein de las
dobles trasposiciones (\omterm{def:b3:groups:normal}{normal} en $S_4$:
es una unión de \omterm{ex:b3:groups:actions}{clases de conjugación}),
con cocientes abelianos $\Z/2\Z$, $\Z/3\Z$, $V$. En
el~\cref{ch:b3:galois}, «la ecuación general de grado $n$ es
\omterm{def:b3:groups:derived}{resoluble} por radicales»
\emph{significará} literalmente «$S_n$ es un
\omterm{def:b3:groups:derived}{grupo resoluble}». De ahí la importancia
de la definición siguiente.
\end{example}

\begin{definition}\label{def:b3:groups:simple}
Un grupo $G \neq \{e\}$ es \emph{simple}\index{grupo simple} si sus
únicos \omterm{def:b3:groups:normal}{subgrupos normales} son $\{e\}$ y
$G$. Un grupo simple no abeliano no es
\omterm{def:b3:groups:derived}{resoluble}: $D(G) \trianglelefteq G$ no
es $\{e\}$ (si lo fuera, $G$ sería abeliano), luego $D(G) = G$ y la
\omterm{def:b3:groups:derived}{serie derivada} es constante. Los grupos
simples abelianos son exactamente los $\Z/p\Z$ con $p$ primo (un grupo
abeliano es simple si y solo si no tiene ningún subgrupo propio no
trivial, si y solo si es cíclico de orden primo, por Lagrange).
\end{definition}

\begin{lemma}\label{lem:b3:groups:threecycles}
Para $n \geq 3$, $A_n$ está generado por los $3$-ciclos; para
$n \geq 5$, todos los $3$-ciclos son conjugados \emph{en $A_n$}.
\end{lemma}

\begin{proof}
Un elemento de $A_n$ es producto de un número par de trasposiciones;
agrupémoslas de dos en dos y usemos (la composición se lee de derecha a
izquierda)
\[
(a\,b)(c\,d) = (a\,c\,b)(a\,c\,d), \qquad
(a\,b)(b\,c) = (a\,b\,c), \qquad
(a\,b)(a\,b) = e
\]
para pares disjuntos, solapados e iguales respectivamente: cada par de
trasposiciones es un producto de $3$-ciclos.

Conjugación: $\sigma(a\,b\,c)\sigma^{-1} = (\sigma a\; \sigma b\;
\sigma c)$, luego dos $3$-ciclos cualesquiera son conjugados por algún
$\sigma \in
S_n$. Si $\sigma$ es impar, sustitúyase por $\sigma' = \sigma (d\,e)$,
donde $d, e$ son dos puntos fuera de $\{a, b, c\}$ --- existen porque
$n \geq 5$ ---; entonces $\sigma'$ es par y $\sigma'(a\,b\,c)
\sigma'^{-1} = \sigma(a\,b\,c)\sigma^{-1}$, ya que $(d\,e)$ conmuta con
$(a\,b\,c)$.
\end{proof}

\begin{theorem}[Simplicidad del grupo
alternado]\label{thm:b3:groups:ansimple}
$A_n$ es \omterm{def:b3:groups:simple}{simple} para
$n \geq 5$.\index{grupo alternado}
\end{theorem}

\begin{proof}
Sea $N \trianglelefteq A_n$, $N \neq \{e\}$. Por
el~\cref{lem:b3:groups:threecycles} basta ver que $N$ contiene \emph{un}
$3$-ciclo: la normalidad y la conjugación de los $3$-ciclos en $A_n$
ponen entonces todos los $3$-ciclos en $N$, luego $N = A_n$.

Para $\rho \in S_n$, sea $F(\rho) = \{x : \rho(x) \neq x\}$ su
\emph{soporte} y $f(\rho) = \abs{F(\rho)}$. Elijamos $\sigma \in N
\setminus \{e\}$ con $f(\sigma)$ \emph{mínimo}. Obsérvese que una
permutación par no trivial cumple $f \geq 3$ y que $f(\sigma) =
4$ es imposible para $\sigma \in A_n$ salvo que $\sigma$ sea una doble
trasposición (un $4$-ciclo es impar). Veamos que $\sigma$ es un
$3$-ciclo.

\emph{Caso A: $\sigma$ es un producto de trasposiciones disjuntas},
digamos $\sigma = (a\,b)(c\,d)\cdots$ con $f(\sigma) \geq 4$. Tomemos
$e'
\notin \{a, b, c, d\}$ (es posible: $n \geq 5$), pongamos $\tau =
(c\,d\,e')$ y
\[
\sigma' = \tau\sigma\tau^{-1}\,\sigma^{-1} \in N
\qquad (\tau\sigma\tau^{-1} \in N \text{ por normalidad}).
\]
Como $\sigma\tau^{-1}\sigma^{-1} = (\sigma c\;\sigma e'\;\sigma
d) = (d\;\sigma e'\;c)$ (usando $\sigma c = d$, $\sigma d = c$),
obtenemos $\sigma' = (c\,d\,e')(d\;\sigma e'\;c)$.

Si $\sigma e' = e'$ (lo que ocurre en particular cuando $f(\sigma) =
4$, es decir, $\sigma = (a\,b)(c\,d)$): entonces $(d\,e'\,c) =
(c\,d\,e')$ y $\sigma' = (c\,d\,e')^2 = (c\,e'\,d)$, un $3$-ciclo que
está en $N$, con $f(\sigma') = 3 < 4 \leq f(\sigma)$ --- en
contradicción con la minimalidad.

Si $\sigma e' \neq e'$: entonces $\sigma e' \notin \{a, b, c, d,
e'\}$ ($\sigma$ intercambia $a,b$ y $c,d$, y $e' \notin \{a,b,c,d\}$ con
$\sigma$ inyectiva), luego $\sigma$ mueve los seis puntos $a, b,
c, d, e', \sigma e'$: $f(\sigma) \geq 6$. Por otra parte $\sigma'$,
producto de dos $3$-ciclos con soportes contenidos en $\{c, d,
e', \sigma e'\}$, cumple $f(\sigma') \leq 4$; y $\sigma' \neq
e$, pues $\sigma'(d) = \tau\sigma\tau^{-1}(c) = \tau\sigma(e') =
\sigma e' \neq d$ ($\tau$ fija $\sigma e' \notin \{c,d,e'\}$). Así que
$\sigma' \in N \setminus\{e\}$ con $f(\sigma') \leq 4 <
f(\sigma)$: se contradice la minimalidad.

\emph{Caso B: algún ciclo de $\sigma$ tiene longitud $\geq 3$}, digamos
$\sigma(a) = b$, $\sigma(b) = c$ con $a, b, c$ distintos. Si $\sigma$ es
exactamente ese $3$-ciclo, hemos terminado. En caso contrario
$f(\sigma) \geq 5$ (el caso $f(\sigma) = 4$ con un $\geq3$-ciclo daría
el $4$-ciclo impar, excluido), de modo que podemos tomar $d, e' \in
F(\sigma) \setminus \{a, b, c\}$. Pongamos $\tau = (c\,d\,e')$ y
$\sigma' = \tau\sigma\tau^{-1}\sigma^{-1} \in N$. Como antes,
$\sigma' = (c\,d\,e')\,(\sigma c\;\sigma e'\;\sigma d)$ solo mueve
puntos de
\[
M = \{c, d, e'\} \cup \{\sigma c, \sigma d, \sigma e'\}
\subseteq F(\sigma)
\]
(las imágenes de los puntos movidos se mueven: $\sigma(x) \ne x$ implica
$\sigma(\sigma x) \neq \sigma x$, por ser $\sigma$ inyectiva). Además
$b \notin M$: los cinco puntos $a, b, c, d, e'$ son distintos, luego $b
\notin \{c, d, e'\}$; y $b \in \{\sigma c, \sigma d, \sigma
e'\}$ obligaría a $a \in \{c, d, e'\}$ (aplíquese $\sigma^{-1}$, usando
$\sigma a = b$), lo cual es falso. Por tanto $\sigma'$ fija $b$,
mientras que $\sigma$ mueve $b$; y $F(\sigma') \subseteq F(\sigma)$.
Finalmente $\sigma' \neq e$:
$\sigma^{-1}(c) = b$, $\tau^{-1}(b) = b$, $\sigma(b) = c$,
$\tau(c) = d$, luego $\sigma'(c) = d \neq c$. Así $\sigma' \in N
\setminus \{e\}$ con $f(\sigma') \leq f(\sigma) - 1$, en contradicción
con la minimalidad.

Siendo imposibles ambos casos, $\sigma$ es un $3$-ciclo.
\end{proof}

\begin{corollary}\label{cor:b3:groups:snnotsolvable}
Para $n \geq 5$: $A_n$ y $S_n$ no son
\omterm{def:b3:groups:derived}{resolubles}, y los únicos
\omterm{def:b3:groups:normal}{subgrupos normales} de $S_n$ son $\{e\}$,
$A_n$ y $S_n$.
\end{corollary}

\begin{proof}
$A_n$ es \omterm{def:b3:groups:simple}{simple} no abeliano, luego no es
\omterm{def:b3:groups:derived}{resoluble} (\cref{def:b3:groups:simple});
y un grupo que contiene un subgrupo no
\omterm{def:b3:groups:derived}{resoluble} no es \omterm{def:b3:groups:derived}{resoluble}
(la~\cref{prop:b3:groups:derived}). Sea $N
\trianglelefteq S_n$: entonces $N \cap A_n \trianglelefteq A_n$ es igual
a $\{e\}$ o a $A_n$. Si $N \cap A_n = A_n$, entonces $A_n \subseteq N$ y
$N \in \{A_n, S_n\}$ por el índice. Si $N \cap A_n = \{e\}$, la
restricción a $N$ de la proyección $S_n \to S_n/A_n \cong
\Z/2\Z$ es inyectiva, luego $\abs N \leq 2$; si $N = \{e, \sigma\}$, la
normalidad hace que la \omterm{ex:b3:groups:actions}{clase de
conjugación} de $\sigma$ sea igual a $\{\sigma\}$, es decir,
$\sigma \in Z(S_n)$. Pero $Z(S_n) = \{e\}$ para $n \geq 3$: si
$\sigma \ne e$ lleva $a$ a $b \neq a$, tómese $c
\notin \{a, b\}$; entonces $(b\,c)\sigma(b\,c)^{-1}$ lleva $a$ a $c
\ne b$, luego difiere de $\sigma$. Por tanto $N = \{e\}$.
\end{proof}

\begin{theorem}[Jordan--Hölder]\label{thm:b3:groups:jordanholder}
Todo grupo finito $G \neq \{e\}$ admite una \emph{serie de
composición}\index{serie de composición}
\[
\{e\} = G_0 \trianglelefteq G_1 \trianglelefteq \cdots
\trianglelefteq G_r = G,
\qquad G_{i}/G_{i-1} \text{ simple},
\]
y el multiconjunto de los \emph{factores de composición} $G_i/G_{i-1}$,
salvo isomorfismo, no depende de la serie elegida. Un grupo finito es
\omterm{def:b3:groups:derived}{resoluble} si y solo si todos sus
factores de composición son cíclicos de orden primo.\index{teorema de
Jordan--Hölder}
\end{theorem}

\begin{proof}
\emph{Existencia}: por inducción sobre $\abs G$. Si $G$ es
\omterm{def:b3:groups:simple}{simple}, tómese $\{e\} \trianglelefteq G$.
En otro caso, elíjase un \omterm{def:b3:groups:normal}{subgrupo normal}
propio maximal $N$ (solo hay un número finito de subgrupos); $G/N$ es
\omterm{def:b3:groups:simple}{simple} por el teorema de correspondencia
(un \omterm{def:b3:groups:normal}{subgrupo normal} propio y no trivial
de $G/N$ se levantaría a un \omterm{def:b3:groups:normal}{subgrupo
normal} de $G$ estrictamente comprendido entre $N$ y $G$). Añádase
$N \trianglelefteq G$ a una \omterm{thm:b3:groups:jordanholder}{serie de
composición} de $N$.

\emph{Unicidad}: por inducción sobre $\abs G$, siendo claro el caso $G$
\omterm{def:b3:groups:simple}{simple}. Tomemos dos
\omterm{thm:b3:groups:jordanholder}{series de composición}, de
penúltimos términos $M \trianglelefteq G$ y $N \trianglelefteq G$ (de
modo que $G/M$, $G/N$ son \omterm{def:b3:groups:simple}{simples}). Si
$M = N$, se concluye por inducción aplicada a $M$. En otro caso $MN$,
\omterm{def:b3:groups:normal}{normal} en $G$ y que contiene
estrictamente a $M$, es igual a $G$ ($M$ es
\omterm{def:b3:groups:normal}{normal} maximal: todo
\omterm{def:b3:groups:normal}{normal} $M \subsetneq L
\subsetneq G$ se aplicaría en un \omterm{def:b3:groups:normal}{subgrupo
normal} propio y no trivial del \omterm{def:b3:groups:simple}{simple}
$G/M$). El segundo teorema de isomorfía da
\[
G/M = MN/M \cong N/(M \cap N),
\qquad
G/N = MN/N \cong M/(M \cap N).
\]
Pongamos $K = M \cap N$ ($\trianglelefteq G$) y fijemos una
\omterm{thm:b3:groups:jordanholder}{serie de composición} de $K$.
Entonces $M$ admite dos \omterm{thm:b3:groups:jordanholder}{series de
composición}: la original y la de $K$ seguida de $K \trianglelefteq
M$ (el cociente $M/K \cong G/N$ es
\omterm{def:b3:groups:simple}{simple}). Por inducción (aplicada a $M$),
los factores de la serie original de $M$ son
$\{\text{factores de } K\} \cup \{G/N\}$; y otro tanto para $N$. Por
tanto, ambas series de $G$ tienen por factores
\[
\{\text{factores de } K\} \;\cup\; \{\,G/N,\; G/M\,\},
\]
el mismo multiconjunto.

\omterm{def:b3:groups:derived}{Resolubilidad}: si todos los factores son
$\Z/p_i\Z$, la serie es una cadena de cocientes abelianos, luego $G$ es
\omterm{def:b3:groups:derived}{resoluble}
(la~\cref{prop:b3:groups:derived}). Recíprocamente, un factor de
composición de un \omterm{def:b3:groups:derived}{grupo resoluble} es
\omterm{def:b3:groups:derived}{resoluble} (cociente de un subgrupo) y
\omterm{def:b3:groups:simple}{simple}; y un
\omterm{def:b3:groups:simple}{grupo simple}
\omterm{def:b3:groups:derived}{resoluble} es abeliano ($D(G) \ne G$
obliga a $D(G) = \{e\}$), luego es algún $\Z/p\Z$.
\end{proof}

\begin{remark}\label{rem:b3:groups:jhmeaning}
Jordan--Hölder dice que todo grupo finito está construido a partir de
\omterm{def:b3:groups:simple}{grupos simples}, con una lista de piezas
bien definida --- una aritmética de grupos en la que los
\omterm{def:b3:groups:simple}{grupos simples} son los primos y donde
\emph{la manera de pegar las piezas} (los datos de extensión, como en el
\omterm{def:b3:groups:semidirect}{producto semidirecto}) sustituye a la
mera multiplicación. La clasificación de los
\omterm{def:b3:groups:simple}{grupos simples} finitos --- los cíclicos
$\Z/p\Z$, los alternados $A_{n \geq
5}$, dieciséis familias de tipo Lie y $26$ grupos esporádicos --- es uno
de los monumentos de las matemáticas del siglo XX; su demostración,
repartida en unas diez mil páginas de revista, queda muy lejos del
alcance de este curso.
\end{remark}

\begin{omfigure}
\begin{tikzpicture}[xscale=1.55, yscale=1.4]
  \tikzset{sg/.style={draw, rounded corners=2pt, inner sep=3pt,
      fill=omDef!6, font=\small}}
  \node[sg, fill=omThm!10] (G)   at (0, 3)    {$D_4$};
  \node[sg] (K1)  at (-1.6, 2)   {$\langle r^2, s\rangle$};
  \node[sg] (R)   at (0, 2)      {$\langle r\rangle$};
  \node[sg] (K2)  at (1.6, 2)    {$\langle r^2, rs\rangle$};
  \node[sg] (S1)  at (-2.4, 1)   {$\langle s\rangle$};
  \node[sg] (S2)  at (-1.2, 1)   {$\langle r^2 s\rangle$};
  \node[sg, fill=omProp!12] (Z) at (0, 1) {$\langle r^2\rangle$};
  \node[sg] (S3)  at (1.2, 1)    {$\langle rs\rangle$};
  \node[sg] (S4)  at (2.4, 1)    {$\langle r^3 s\rangle$};
  \node[sg] (E)   at (0, 0)      {$\{e\}$};
  \draw (G) -- (K1); \draw (G) -- (R); \draw (G) -- (K2);
  \draw (K1) -- (S1); \draw (K1) -- (S2); \draw (K1) -- (Z);
  \draw (R) -- (Z);
  \draw (K2) -- (Z); \draw (K2) -- (S3); \draw (K2) -- (S4);
  \draw (S1) -- (E); \draw (S2) -- (E); \draw (Z) -- (E);
  \draw (S3) -- (E); \draw (S4) -- (E);
\end{tikzpicture}

{\small Los diez subgrupos del
\omterm{ex:b3:groups:semidirectexamples}{grupo diédrico}
$D_4 = \langle r,
s \mid r^4 = s^2 = e,\ srs^{-1} = r^{-1}\rangle$. Los tres subgrupos de
índice $2$ (fila central) son \omterm{def:b3:groups:normal}{normales},
al igual que el \omterm{ex:b3:groups:actions}{centro}
$\langle r^2\rangle$ (resaltado); los cuatro subgrupos de reflexión se
reparten en dos \omterm{ex:b3:groups:actions}{clases de conjugación} de
dos elementos. Las cadenas de abajo arriba dan
\omterm{thm:b3:groups:jordanholder}{series de composición}, por ejemplo
$\{e\} \trianglelefteq \langle
r^2\rangle \trianglelefteq \langle r\rangle \trianglelefteq D_4$:
factores $\Z/2\Z, \Z/2\Z, \Z/2\Z$ --- siempre el mismo multiconjunto,
como exige Jordan--Hölder.}
\end{omfigure}

\section{Ejercicios}

\begin{exercise}[$\star$]\label{exo:b3:groups:1}
(a) Demostrar que todo subgrupo de índice $2$ es
\omterm{def:b3:groups:normal}{normal}. (b) Demostrar que si
$[G : H] = 2$, entonces $x^2 \in H$ para todo $x \in
G$. (c) Deducir que $A_4$ no tiene ningún subgrupo de orden $6$: el
recíproco de Lagrange es falso. \emph{(Contar los cuadrados de los
$3$-ciclos.)}
\end{exercise}

\begin{exercise}[$\star$]\label{exo:b3:groups:2}
Demostrar que si $G/Z(G)$ es cíclico, entonces $G$ es abeliano. Deducir
de nuevo que todo grupo de orden $p^2$ es abeliano y exhibir, para cada
primo $p$, un grupo no abeliano de orden $p^3$. \emph{(Piénsese en las
matrices triangulares superiores con diagonal unidad sobre
$\mathbb F_p$.)}
\end{exercise}

\begin{exercise}[$\star$]\label{exo:b3:groups:3}
(a) Demostrar que $\operatorname{Aut}(\Z/n\Z) \cong (\Z/n\Z)^\times$.
(b) Demostrar que los automorfismos interiores $\iota_g \colon x \mapsto
gxg^{-1}$ forman un \omterm{def:b3:groups:normal}{subgrupo normal}
$\operatorname{Inn}(G)
\trianglelefteq \operatorname{Aut}(G)$, con
$\operatorname{Inn}(G) \cong G/Z(G)$.
\end{exercise}

\begin{exercise}[$\star\star$]\label{exo:b3:groups:4}
Sean $H, K$ subgrupos de un grupo finito $G$. (a) Demostrar la
\emph{fórmula del producto} $\abs{HK}\,\abs{H\cap K} =
\abs H\, \abs K$ contando las fibras de la aplicación $H \times K
\to HK$, $(h,k) \mapsto hk$. (b) Demostrar que $HK$ es un subgrupo si y
solo si $HK = KH$ (lo cual es automático cuando uno de los dos es
\omterm{def:b3:groups:normal}{normal}). (c) Si $H, K \trianglelefteq G$
y $H \cap K = \{e\}$, demostrar que $hk = kh$ para todos $h \in H$,
$k \in K$.
\end{exercise}

\begin{exercise}[$\star\star$]\label{exo:b3:groups:5}
(Lema de recuento de Burnside) Un grupo finito $G$ actúa sobre un
conjunto finito $X$. Demostrar que el número de
\omterm{def:b3:groups:action}{órbitas} es el \emph{número medio de
puntos fijos}:
\[
\#\{\text{órbitas}\} = \frac{1}{\abs G}\sum_{g \in G}
\abs{\operatorname{Fix}(g)},
\qquad \operatorname{Fix}(g) = \{x \in X : g \cdot x = x\},
\]
contando de dos maneras el conjunto $\{(g,x) : g\cdot x = x\}$.
Aplicación: $\Z/p\Z$ ($p$ primo) actúa por rotación sobre los collares
de $p$ cuentas con $a$ colores disponibles; deducir el pequeño teorema
de Fermat $a^p \equiv a \pmod p$.
\end{exercise}

\begin{exercise}[$\star$]\label{exo:b3:groups:6}
Usando los teoremas de Sylow, demostrar que todo grupo de orden $15$ es
cíclico y que todo grupo de orden $45$ es abeliano.
\end{exercise}

\begin{exercise}[$\star\star$]\label{exo:b3:groups:7}
Demostrar que ningún grupo de orden $30$, ni de orden $56$, es
\omterm{def:b3:groups:simple}{simple}. \emph{(Para $30$: si $n_3 \neq 1$
y $n_5 \neq 1$, contar los elementos de órdenes $3$ y $5$. Para $56$:
contar los elementos de orden $7$.)}
\end{exercise}

\begin{exercise}[$\star\star$]\label{exo:b3:groups:8}
(a) Sea $H \leq G$ de índice $n$. Demostrar que la
\omterm{def:b3:groups:action}{acción} de $G$ sobre $G/H$ da un morfismo
$G \to S_n$ cuyo núcleo $\bigcap_{g \in
G} gHg^{-1}$ es el mayor \omterm{def:b3:groups:normal}{subgrupo normal}
de $G$ contenido en
$H$.
(b) Deducir que, si $G$ es finito y $p$ es el \emph{menor} divisor primo
de $\abs G$, todo subgrupo de índice $p$ es
\omterm{def:b3:groups:normal}{normal}. (c) Demostrar que toda
\omterm{def:b3:groups:action}{acción} transitiva de $G$ sobre un
conjunto $X$ es isomorfa a la \omterm{def:b3:groups:action}{acción}
sobre un espacio de clases: existe una biyección $X
\to G/G_x$ que conmuta con las acciones.
\end{exercise}

\begin{exercise}[$\star\star$]\label{exo:b3:groups:9}
(a) Demostrar que $D(G)$ es el menor
\omterm{def:b3:groups:normal}{subgrupo normal} de $G$ con cociente
abeliano, y que todo morfismo de $G$ en un grupo abeliano se factoriza
de manera única a través de la \emph{abelianización}
$G^{\mathrm{ab}} = G/D(G)$.
(b) Calcular $D(S_n)$ y $S_n^{\mathrm{ab}}$ para $n \geq 2$, así como
$D(Q_8)$ y $Q_8^{\mathrm{ab}}$.
\end{exercise}

\begin{exercise}[$\star\star$]\label{exo:b3:groups:10}
Sea $G$ un $p$-grupo y $H \subsetneq G$ un subgrupo propio. Demostrar
que $H \subsetneq N_G(H)$ («los
\omterm{ex:b3:groups:actions}{normalizadores} crecen») y deducir que
todo subgrupo maximal de un $p$-grupo es
\omterm{def:b3:groups:normal}{normal} de índice $p$. \emph{(Inducción
sobre $\abs G$, usando $Z(G) \neq \{e\}$: trátense por separado los
casos $Z(G) \subseteq H$ y $Z(G) \not\subseteq H$.)}
\end{exercise}

\begin{exercise}[$\star\star\star$]\label{exo:b3:groups:11}
(Simplicidad de $A_5$, a mano) (a) Demostrar que las
\omterm{ex:b3:groups:actions}{clases de conjugación} de $A_5$ tienen
cardinales $1$, $15$, $20$, $12$, $12$. Préstese atención al
desdoblamiento de la $S_5$-clase de los $5$-ciclos: para un $5$-ciclo
$\sigma$, compárense los \omterm{ex:b3:groups:actions}{centralizadores}
de $\sigma$ en $S_5$ y en $A_5$. (b) Deducir que $A_5$ es
\omterm{def:b3:groups:simple}{simple}: un
\omterm{def:b3:groups:normal}{subgrupo normal} es una unión de
\omterm{ex:b3:groups:actions}{clases de conjugación}, contiene $e$ y su
cardinal divide a
$60$.
(c) Demostrar que un \omterm{def:b3:groups:simple}{grupo simple} de
orden $60$ cumple necesariamente $n_5 =
6$.
\end{exercise}

\begin{exercise}[$\star\star$]\label{exo:b3:groups:12}
(Los \omterm{ex:b3:groups:actions}{normalizadores} de los
\omterm{def:b3:groups:sylow}{subgrupos de Sylow} son autonormalizantes)
Sea $P$ un $p$-subgrupo de Sylow de un grupo finito $G$ y sea $H =
N_G(P)$. (a) Demostrar que $P$ es el \emph{único} $p$-subgrupo de Sylow
de
$H$.
(b) Deducir $N_G(H) = H$. \emph{(Para $g \in N_G(H)$: $gPg^{-1}$ es un
$p$-subgrupo de Sylow de $H$, luego $gPg^{-1} = P$.)} (c) Concluir que
ningún \omterm{ex:b3:groups:actions}{normalizador} de Sylow está
contenido en un \omterm{def:b3:groups:normal}{subgrupo normal} propio de
$G$, y que todo subgrupo maximal que contenga a $N_G(P)$ es
autonormalizante.
\end{exercise}

\section{Problema: los grupos de orden a lo sumo 15}

\begin{problem}[{Problema de fin de semana --- clasificación de los
grupos
pequeños}]\label{pb:b3:groups:1} El objetivo es una clasificación
completa, con demostraciones detalladas, de los grupos de orden
$\leq 15$ salvo isomorfismo. Los órdenes $1, 2, 3, 5,
7, 11, 13$ los resuelve Lagrange (son cíclicos), y los órdenes $4$ y $9$
el~\cref{thm:b3:groups:pfixed} junto con el análisis de $p^2$ que sigue:
quedan $6, 8, 10, 12, 14, 15$.

\medskip
\noindent\textbf{Parte I --- Herramientas.}
\begin{enumerate}
  \item Demostrar que un grupo en el que todo elemento cumple $x^2 =
        e$ es abeliano; deducir que un grupo finito de este tipo tiene
  orden $2^k$ y es isomorfo a $(\Z/2\Z)^k$. \emph{(Véase como espacio
  vectorial sobre $\mathbb F_2$.)} \item Demostrar que un grupo de orden
  $p^2$ es isomorfo a $\Z/p^2\Z$ o a $(\Z/p\Z)^2$. Listar los grupos
  abelianos de orden $8$ salvo isomorfismo: $\Z/8\Z$, $\Z/4\Z \times
        \Z/2\Z$, $(\Z/2\Z)^3$ --- probar que la lista es completa y sin
  repeticiones \emph{sin} recurrir al teorema de estructura
  del~\cref{ch:b3:modules} (discútase según el orden máximo de un
  elemento). \item Sean
  $\varphi, \varphi' \colon K \to \operatorname{Aut}(H)$ dos
  acciones. Demostrar que si
  $\varphi' = \varphi \circ
        \alpha$ con $\alpha \in \operatorname{Aut}(K)$, entonces $H
        \rtimes_{\varphi} K \cong H \rtimes_{\varphi'} K$. \item
  Determinar $\operatorname{Aut}(\Z/n\Z)$ explícitamente para
  $n = 3, 4, 5,
        7$ y demostrar que $\operatorname{Aut}\bigl((\Z/2\Z)^2
        \bigr) \cong S_3$.
\end{enumerate}

\medskip
\noindent\textbf{Parte II --- Órdenes $2p$ ($6$, $10$, $14$) y
$pq$.}
\begin{enumerate}[resume]
  \item Sea $\abs G = 2p$ con $p$ primo impar. Demostrar que $G$ tiene
  un \omterm{def:b3:groups:normal}{subgrupo normal}
  $N = \langle r \rangle$ de orden $p$ y un elemento $s$ de orden $2$
  fuera de $N$. \item Deducir que
  $G \cong \Z/p\Z \rtimes_\varphi \Z/2\Z$, donde
  $\varphi(1) \in \operatorname{Aut}(\Z/p\Z)$ es una involución, y
  concluir: $G \cong \Z/2p\Z$ o $G \cong
        D_p$; comprobar que estos dos no son isomorfos. Esto resuelve
  los órdenes $6$, $10$ y $14$. \item Más en general, sea $\abs G = pq$
  con $p < q$ primos. Demostrar que si $p \nmid q-1$ entonces $G$ es
  cíclico (\cref{ex:b3:groups:pq}), y que si $p \mid q - 1$ existe,
  además de $\Z/pq\Z$, \emph{exactamente un} grupo no abeliano
  $\Z/q\Z \rtimes \Z/p\Z$ salvo isomorfismo --- utilícese la pregunta 3
  y el hecho de que $\operatorname{Aut}(\Z/q\Z) \cong
        (\Z/q\Z)^\times$ es cíclico de orden $q - 1$, admitido aquí y
  demostrado en el~\cref{ch:b3:galois} (ciclicidad de
  $\mathbb F_q^\times$). Concluir para el orden $15$.
\end{enumerate}

\medskip
\noindent\textbf{Parte III --- Orden $8$.} Sea $G$ no abeliano de orden
$8$.
\begin{enumerate}[resume]
  \item Demostrar que $G$ tiene un elemento $r$ de orden $4$ (úsese la
  pregunta 1) y que $N = \langle r\rangle$ es
  \omterm{def:b3:groups:normal}{normal}. \item Sea $s \notin N$.
  Demostrar que $s^2 \in N$ (\cref{exo:b3:groups:1}(b)), que
  $srs^{-1} = r^{-1}$ (examínense las posibles imágenes de $r$ por
  conjugación, que han de tener orden 4, y descártese $srs^{-1} = r$) y
  que $s^2 \in \{e, r^2\}$ \emph{(¿qué ocurre si $s^2 = r$ o $r^3$?, ¿y
  por qué ha de conmutar $s^2$ con $s$?)}. \item En el caso $s^2 = e$,
  demostrar que $G \cong D_4$. \item En el caso $s^2 = r^2$, demostrar
  que la tabla de multiplicar queda enteramente determinada; el grupo
  resultante es el \emph{grupo de cuaterniones}\index{grupo de
  cuaterniones} $Q_8 =
        \{\pm 1, \pm \mathrm i, \pm \mathrm j, \pm \mathrm k\}$,
  $\mathrm i^2 = \mathrm j^2 = \mathrm k^2 = \mathrm i\,
        \mathrm j\,\mathrm k = -1$ (póngase $r = \mathrm i$, $s =
        \mathrm j$). Comprobar que $Q_8$ existe, por ejemplo dentro de
  $GL_2(\C)$ mediante
        \[
        \mathrm i \mapsto \begin{pmatrix} \iu & 0\\ 0 & -\iu
        \end{pmatrix},
        \qquad
        \mathrm j \mapsto \begin{pmatrix} 0 & 1\\ -1 & 0
        \end{pmatrix}.
        \]
  \item Demostrar que todo subgrupo no trivial de $Q_8$ contiene $-1$;
  deducir que todo subgrupo de $Q_8$ es
  \omterm{def:b3:groups:normal}{normal}, que $D_4 \not\cong Q_8$
  (cuéntense los elementos de orden $2$) y que $Q_8$ no es
  \omterm{def:b3:groups:semidirect}{producto semidirecto} de dos
  subgrupos propios.
\end{enumerate}

\medskip
\noindent\textbf{Parte IV --- Orden $12$.} Sean $\abs G = 12$, $P_3
\in \mathrm{Syl}_3(G)$, $P_2 \in \mathrm{Syl}_2(G)$.
\begin{enumerate}[resume]
  \item Demostrar que $n_3 \in \{1, 4\}$, $n_2 \in \{1, 3\}$, y que $n_3
        = 4$ obliga a $n_2 = 1$ \emph{(cuéntense los elementos de orden
        $3$)}.
  \item Supóngase $n_3 = 4$. La \omterm{def:b3:groups:action}{acción}
  por conjugación sobre $\mathrm{Syl}_3$ da $\rho \colon G \to S_4$.
  Demostrar que $\ker \rho$, contenido en todos los $N_G(P_3)$ y por
  tanto de orden divisor de $3$, es trivial \emph{(¿por qué no puede
  tener orden $3$?)}; que la imagen, subgrupo de orden $12$ de $S_4$, es
  necesariamente $A_4$ \emph{(los subgrupos de índice 2 son
  \omterm{def:b3:groups:normal}{normales} y contienen todos los
  cuadrados --- el~\cref{exo:b3:groups:1}; cuéntense los cuadrados de
  $S_4$)}; y concluir que $G \cong A_4$. \item Supóngase $n_3 = 1$, de
  modo que $G \cong \Z/3\Z \rtimes_\varphi P_2$ con
  $\varphi \colon P_2 \to \operatorname{Aut}(\Z/3\Z)
        \cong \Z/2\Z$. Enumérense los casos: $\varphi$ trivial da
  $\Z/12\Z$ y $\Z/6\Z \times \Z/2\Z$; $P_2 =
        \Z/4\Z$ con $\varphi$ sobreyectivo da el \emph{grupo dicíclico}
  $\mathrm{Dic}_3 = \Z/3\Z \rtimes
        \Z/4\Z$; $P_2 = (\Z/2\Z)^2$ con $\varphi$ sobreyectivo da, salvo
  la equivalencia de la pregunta 3, un único grupo --- demuéstrese que
  es $D_6$, por ejemplo exhibiendo un elemento de orden $6$ y una
  involución de tipo reflexión. \item Demostrar que $\Z/12\Z$,
  $\Z/6\Z \times \Z/2\Z$, $D_6$, $A_4$, $\mathrm{Dic}_3$ son dos a dos
  no isomorfos \emph{(cuéntense los elementos de orden $2$, o úsese
  $n_3$)}. Esto resuelve el orden
        $12$.
\end{enumerate}

\medskip
\noindent\textbf{Parte V --- Síntesis.}
\begin{enumerate}[resume]
  \item Montar la tabla de clasificación: para cada orden $n \leq
        15$, la lista completa de los grupos salvo isomorfismo, con los
  recuentos $1, 1, 1, 2, 1, 2, 1, 5, 2, 2, 1, 5, 1, 2, 1$.
\end{enumerate}

\medskip
\noindent\textbf{Parte VI --- Más allá: los grupos de orden $p^3$ con
$p$ impar.} El análisis del orden $8$ de la Parte III tiene un bello
análogo para primos impares, con un fenómeno genuinamente nuevo. Sea $p$
un primo impar y sea $G$ no abeliano de orden $p^3$.
\begin{enumerate}[resume]
  \item Demostrar que $\abs{Z(G)} = p$, que $G/Z(G) \cong
        (\Z/p\Z)^2$ \emph{(un cociente cíclico por el
  \omterm{ex:b3:groups:actions}{centro} obliga a la abelianidad:
  el~\cref{exo:b3:groups:2})} y que $D(G) =
        Z(G)$ \emph{(para $D(G) \subseteq Z(G)$, úsese que $G/Z(G)$ es
  abeliano; para la igualdad, $G$ es no abeliano y $D(G) \neq \{e\}$)}.
  Deducir que todo \omterm{def:b3:groups:derived}{conmutador}
  $[x, y] = xyx^{-1}y^{-1}$ es central y de orden divisor de $p$. \item
  (La identidad clave) Sean $x, y \in G$ y $z = [y, x]$, central.
  Demostrar por inducción sobre $k$:
        \[
        (xy)^k = x^k y^k z^{k(k-1)/2} .
        \]
        \emph{(Pásese cada $y$ por delante de cada $x$; cada cruce
        cuesta un factor central $z$.)} \item Deducir que, para $p$
        impar, la aplicación $\theta\colon x \mapsto
        x^p$ es un morfismo de grupos de $G$ en $Z(G)$ \emph{(¿por qué
        es $x^p$ central?, ¿por qué $z^{p(p-1)/2} = e$ necesita que $p$
        sea impar?)}, y concluir que $G$ tiene exponente $p$ o $p^2$,
        distinguiéndose ambos casos según que $\theta$ sea trivial o no.
        \item (Exponente $p$) Supóngase que todo elemento cumple
        $x^p = e$. Tómense $x, y$ cuyas clases generen $G/Z(G)$ y
        póngase $z =
        [y, x]$. Demostrar que $z \neq e$, que todo elemento de $G$ se
        escribe de manera única como $x^ay^bz^c$ ($0 \leq a, b, c < p$)
        y que la multiplicación queda enteramente determinada por las
        relaciones $x^p = y^p = z^p = e$, $z$ central, $[y, x] =
        z$. Comprobar que el \emph{grupo de Heisenberg}
        \[
        H_p = \left\{
        \begin{pmatrix} 1 & a & c\\ 0 & 1 & b\\ 0 & 0 & 1
        \end{pmatrix} : a, b, c \in \mathbb F_p \right\}
        \subseteq GL_3(\mathbb F_p)
        \]
        realiza estas relaciones y tiene exponente $p$ (calcúlese
        $(I + N)^p$ con $N$ estrictamente triangular superior, usando
        $N^3 = 0$ y $p \geq 3$): todo grupo no abeliano de orden $p^3$ y
        exponente $p$ es isomorfo a $H_p$. \item (Exponente $p^2$)
        Supóngase que algún $r \in G$ tiene orden $p^2$ y póngase
        $N = \langle r\rangle$, \omterm{def:b3:groups:normal}{normal}
        (de índice $p$: el~\cref{exo:b3:groups:10}). Demostrar que
        existe $s \notin N$ con $s^p = e$ \emph{(tómese cualquier
        $t \notin N$; usando la pregunta 20, corríjase:
        $\theta(t) = t^p \in Z(G) \subseteq N$ --- justifíquese
        $Z(G) = \langle r^p\rangle$ --- y elíjase $a$ con $s = tr^{a}$
        cumpliendo $s^p = e$; ¿dónde se usa que $p$ es impar?)}.
        Demostrar que $srs^{-1} = r^{1+p}$ salvo sustituir $s$ por una
        potencia suya, y concluir: hay exactamente \emph{un} grupo no
        abeliano de orden $p^3$ y exponente $p^2$, a saber
        $\Z/p^2\Z \rtimes_\varphi \Z/p\Z$ con
        $\varphi(1)\colon r \mapsto r^{1+p}$ \emph{(úsese la pregunta 3;
        $\operatorname{Aut}(\Z/p^2\Z)$ es cíclico de orden $p(p-1)$,
        admitido aquí, luego tiene un único subgrupo de orden $p$)}.
        \item Concluir el recuento: para $p$ impar hay exactamente $5$
        grupos de orden $p^3$ (tres abelianos y dos no abelianos), igual
        que para $p = 2$ --- pero los dos no abelianos ya no son $D_4$ y
        $Q_8$. Localizar con precisión dónde falla el argumento para $p$
        impar cuando $p = 2$: en la identidad de la pregunta 19,
        $z^{k(k-1)/2}$ para $k = p = 2$ vale $z^1 \neq
        e$, de modo que elevar al cuadrado no es un morfismo --- y, en
        efecto, $Q_8$ tiene un único elemento de orden $2$, mientras que
        $D_4$, de exponente $4$, tiene cinco.
\end{enumerate}

\medskip
\noindent\textbf{Parte VII --- Complementos.}
\begin{enumerate}[resume]
  \item Para $p$ impar, contar los elementos de orden $p$ en cada uno de
  los dos grupos no abelianos de orden $p^3$: demostrar que $H_p$ tiene
  exactamente $p^3 - 1$ y que $M_p =
        \Z/p^2\Z \rtimes \Z/p\Z$ tiene exactamente $p^2 - 1$
  \emph{(úsese el morfismo $\theta$ de la pregunta 20: identifíquese su
  imagen y después el orden de su núcleo)}. Comprobarlo numéricamente
  para $p = 3$: $26$ frente a $8$. Explicar por qué ningún argumento de
  este tipo, basado en el morfismo de elevar al cuadrado, puede separar
  $D_4$ de $Q_8$, y qué recuento sí los separa. \item Se dice que un
  entero $n \geq 1$ es \emph{cíclico} si todo grupo de orden $n$ es
  cíclico. Demostrar que si $p^2 \mid n$ para algún primo $p$, o si $n$
  tiene divisores primos $p < q$ con $p \mid q - 1$, entonces $n$ no es
  cíclico \emph{(en cada caso, exhíbase un grupo no cíclico de orden
  $n$, usando la Parte II para el segundo)}. Deducir que $n$ cíclico
  obliga a $\gcd(n, \varphi(n)) = 1$, donde $\varphi$ es la indicatriz
  de Euler, y contrastarlo con la tabla de la pregunta 17: entre
  $n \leq 15$, los órdenes con un único grupo son exactamente
  $n \in \{1, 2, 3, 5, 7, 11, 13, 15\}$, precisamente aquellos con
  $\gcd(n, \varphi(n)) = 1$.
\end{enumerate}
\end{problem}
